\section{Deep Learning para Biomarcadores}
\label{sec:ia-biomarcadores}

O refinamento e barateamento do arcabouço computacional abriram as portas para que a IA atue como uma "lente espectral contínua" no corpo humano. A disciplina de descobrimento de biomarcadores apoiados por IA assume que padrões ininteligíveis ao globo ocular do clínico mais sênior (texturas de frequência alta, perturbações microvasculares no esmalte e osso) são não apenas visíveis, mas altamente correlacionáveis aos perfis moleculares patogênicos e ao prognóstico \cite{shen2024}.

\textbf{Radiômica ontológica e oncologia:} A conversão de tomografias em gigantescas tabelas matriciais repletas de cálculos de correlação textura/heterogeneidade permite ao \textit{deep learning} classificar a periculosidade biológica e invasividade tecidual de cistos odontogênicos ou de carcinomas epidermoides intra-ósseos meses ou anos antes do colapso cortical, ditando ressecções muito mais precisas, seletivas e profiláticas.

\textbf{Monitoramento in loco de osseointegração:} Modelos de detecção computam, via séries longitudinais de radiografias periapicais de controle, a conectividade celular ao titânio (ou zircônia) em estágios imaturos, disparando gatilhos e alertas vermelhos via plataformas em nuvem ao detectar lises ou cavitações inframilimétricas incompatíveis com o ritmo normal de neoformação óssea.

Desta forma, os gêmeos digitais cumprem o pressuposto maior da Odontologia de Precisão (P4): tornam o cuidado participativo, preventivo, preditivo e altamente personalizado. Abandonando a postura falimentar de "tratar o defeito já exaurido", a especialidade evolui para a manutenção assídua do balanço homeostático \cite{technologysociety2025}.
